\section{Products of currents and the BV operator}
\label{sec:hj-bv-products}

A current map gives individual insertions and their higher relations.
Products of insertions require an algebra as well. The Lie space
$S_{N,L}^0$ consists of mixed quadratics, so it cannot contain a
general product of two of its elements. Its vanishing Laplacian
does not justify setting the Laplacian to zero on those products.
A contraction can join two distinct current factors.

Fix $L>0$ throughout the following construction. All degrees are
cohomological degrees of unshifted observables, and $\hbar$ has
degree zero. The BV convention is
\begin{equation}\label{eq:hj-bv-bracket}
\{F,G\}_L=(-1)^{|F|}
 \left(\Delta_L(FG)-(\Delta_LF)G-(-1)^{|F|}F\Delta_LG\right).
\end{equation}
Thus $\Delta_L$ and the bracket have degree one. Products obey
$FG=(-1)^{|F||G|}GF$. The algebra unit has degree zero.

First suppose $F,G,H$ are homogeneous and individually
$\Delta_L$-closed. Write $f=|F|$, $g=|G|$, $h=|H|$.
The binary product identity gives
\begin{equation}\label{eq:hj-bv-pair}
\boxed{\Delta_L(FG)=(-1)^f\{F,G\}_L.}
\end{equation}
The first-slot product rule for the odd bracket is
\[
\{FG,H\}_L
 =F\{G,H\}_L+(-1)^{g(h+1)}\{F,H\}_LG.
\]
Applying \eqref{eq:hj-bv-bracket} to $(FG)H$ therefore gives
\begin{equation}\label{eq:hj-bv-triple}
\boxed{\begin{aligned}
\Delta_L(FGH)
={}&(-1)^f\{F,G\}_LH
   +(-1)^{f+g}F\{G,H\}_L\\
 &+(-1)^{f+gh}\{F,H\}_LG.
\end{aligned}}
\end{equation}
The order of the remaining factors is part of this formula.
If the three pair brackets are also $\Delta_L$-closed, a second
application gives
\begin{equation}\label{eq:hj-bv-triple-square}
\begin{aligned}
\Delta_L^2(FGH)
={}&(-1)^{g+1}\{\{F,G\}_L,H\}_L
   +(-1)^g\{F,\{G,H\}_L\}_L\\
 &+(-1)^{gh+h+1}\{\{F,H\}_L,G\}_L=0.
\end{aligned}
\end{equation}
The final equality is precisely the odd Jacobi identity, with
the displayed factor orders. It verifies the ternary signs
without assuming that products remain in the kernel of $\Delta_L$.

The first internal jet supplies a finite sign calculation.
Take odd variables $\beta_x,\beta_y$ of degree one and even
variables $\gamma_x,\gamma_y$ of degree minus two. Put
\[
\Delta=-\partial_{\beta_x}\partial_{\gamma_x}
       -\partial_{\beta_y}\partial_{\gamma_y},
\qquad
F=\beta_x\gamma_y,\qquad G=\beta_y\gamma_x.
\]
The paired degrees sum to minus one, so $\Delta$ has degree one.
Left differentiation in the odd variables gives
\begin{equation}\label{eq:hj-bv-coordinate}
\Delta F=\Delta G=0,\qquad
\Delta(FG)=\beta_x\gamma_x-\beta_y\gamma_y\ne0.
\end{equation}
The two diagonal terms have opposite Laplacians, so the last
expression is $\Delta$-closed. These are the internal matrices
$e,f,h$ from the first jet. A spatial current calculation below
will exhibit the same nonzero product contraction in the actual
distribution algebra.

\section{The algebra generated by actual current kernels}
\label{sec:hj-bv-algebra}

To retain multiplication, use the current kernels as elements of
their existing observable algebra. This preserves relations that
would disappear if the currents were replaced by independent symbols.
Define the unshifted current space
\begin{equation}\label{eq:hj-bv-current-space}
C_{N,L}=\{F\in\mathcal O_{\mathrm{dist},N}:uF\in S_{N,L}^0\}.
\end{equation}
The preceding current construction proves
\begin{equation}\label{eq:hj-bv-current-closure}
QC_{N,L}\subset C_{N,L},\qquad
\{C_{N,L},C_{N,L}\}_L\subset C_{N,L},\qquad
\Delta_LC_{N,L}=0.
\end{equation}
Every element is a mixed quadratic distribution kernel whose
internal matrix coefficients satisfy \eqref{eq:hj-59}.

Let $\operatorname{Sym}_{\mathbb C}(C_{N,L})$ be the algebraic
graded symmetric algebra, with its unit in symmetric power zero.
Multiplication of actual observables gives an algebra map
\[
\mu_N:\operatorname{Sym}_{\mathbb C}(C_{N,L})
          \longrightarrow\mathcal O_{\mathrm{dist},N}.
\]
Define
\begin{equation}\label{eq:hj-bv-image}
\boxed{A^0_{N,L}=\operatorname{im}\mu_N.}
\end{equation}
An element is a finite sum of finite products of elements of
$C_{N,L}$. The empty product maps to the scalar observable $1$.
The products use external tensor products of distributions,
followed by the graded symmetric coinvariants of their matter slots.
No product of distributions is restricted to a spatial diagonal.

Write $\mathcal I_{N,L}=\ker\mu_N$. Then
\begin{equation}\label{eq:hj-bv-relations}
A^0_{N,L}\cong
\operatorname{Sym}_{\mathbb C}(C_{N,L})/\mathcal I_{N,L}.
\end{equation}
Injectivity of $\mu_N$ is neither needed nor asserted.
All polynomial relations already present among the distribution
observables remain in this quotient.

There is an additional grading by the number of current factors:
\begin{equation}\label{eq:hj-bv-factor-grading}
A^0_{N,L}=\bigoplus_{r\geq0}A^{[r]}_{N,L},\qquad
A^{[r]}_{N,L}=\mu_N\bigl(\operatorname{Sym}^r C_{N,L}\bigr).
\end{equation}
This is not the cohomological grading. Each current factor has
two matter slots, so $A^{[r]}_{N,L}$ lies in observable arity $2r$.
Different arities are distinct direct summands of
$\mathcal O_{\mathrm{dist},N}$. Thus the sum in
\eqref{eq:hj-bv-factor-grading} is direct, and the relation ideal
$\mathcal I_{N,L}$ is homogeneous in this factor grading.
Indeed, if a finite sum of homogeneous components maps to zero,
its components map into distinct observable arities and hence each
maps to zero. Thus every homogeneous component of a relation is
again a relation.

A unital differential BV algebra here is a unital graded commutative
algebra with a degree-one derivation $Q$ and a degree-one operator
$\Delta$ of order at most two, satisfying
$Q^2=\Delta^2=Q\Delta+\Delta Q=0$ and $Q1=\Delta1=0$.
Its bracket is derived by \eqref{eq:hj-bv-bracket}.

\begin{theorem}\label{thm:hj-bv-finite}
The actual image $A^0_{N,L}$ is a unital differential BV subalgebra
of $\mathcal O_{\mathrm{dist},N}$. Its operators are the restrictions
of $Q$ and $\Delta_L$. It is the smallest such subalgebra containing
every unshifted current component.
\end{theorem}
\begin{proof}
The ambient operator $Q$ is a derivation and preserves $C_{N,L}$.
Hence $Q(A^{[r]}_{N,L})\subset A^{[r]}_{N,L}$.
The second-slot product rule is
\[
\{F,GH\}_L=\{F,G\}_LH
 +(-1)^{(|F|+1)|G|}G\{F,H\}_L.
\]
Together with the first-slot rule and
\eqref{eq:hj-bv-current-closure}, it proves
\begin{equation}\label{eq:hj-bv-bracket-grading}
\{A^{[r]}_{N,L},A^{[s]}_{N,L}\}_L
          \subset A^{[r+s-1]}_{N,L}\qquad(r,s\geq1).
\end{equation}
Brackets with the scalar summand vanish.

For homogeneous $F_1,\ldots,F_r\in C_{N,L}$, put
\begin{equation}\label{eq:hj-bv-pair-sign}
\epsilon_{ij}=(-1)^{
 |F_i|\sum_{k<i}|F_k|
 +|F_j|\sum_{\substack{k<j\\k\ne i}}|F_k|}.
\end{equation}
This is the Koszul sign moving $F_i,F_j$, in that order, to the
front. The full product formula is
\begin{equation}\label{eq:hj-bv-all-products}
\boxed{\begin{aligned}
\Delta_L(F_1\cdots F_r)
={}&\sum_{i<j}(-1)^{|F_i|}\epsilon_{ij}\{F_i,F_j\}_L\\
&\hspace{9mm}\cdot
F_1\cdots\widehat F_i\cdots\widehat F_j\cdots F_r.
\end{aligned}}
\end{equation}
The bracket stands first. The other factors retain their original
order, and a hat denotes omission.

For two factors this is \eqref{eq:hj-bv-pair}.
Induct on $r$ by applying \eqref{eq:hj-bv-bracket} to
$F_1(F_2\cdots F_r)$. Terms from the Laplacian of the second
factor account for pairs not containing $1$.
The bracket biderivation accounts for the pairs $(1,j)$.
Moving each chosen pair to the front gives
$\epsilon_{ij}$, and contracting that ordered pair gives
$(-1)^{|F_i|}\{F_i,F_j\}_L$.
These exhaust the terms; single-factor terms vanish by
\eqref{eq:hj-bv-current-closure}. This proves the formula.
In particular,
\begin{equation}\label{eq:hj-bv-lowering}
\Delta_L(A^{[r]}_{N,L})\subset A^{[r-1]}_{N,L},\qquad
\Delta_L(A^{[0]}_{N,L}\oplus A^{[1]}_{N,L})=0.
\end{equation}

These operations are restrictions of the actual ambient operators.
They therefore satisfy
\[
Q^2=\Delta_L^2=Q\Delta_L+\Delta_LQ=0,\qquad
Q1=\Delta_L1=0,
\]
and $\Delta_L$ has differential-operator order at most two.
This order is its order relative to multiplication, not its
cohomological degree or its change of factor grading.

Relations cause no further obstruction. Applying an ambient
operator to a zero distribution gives zero. Equivalently, the
formulas for $Q$, the bracket, and $\Delta_L$ on the symmetric
algebra preserve $\mathcal I_{N,L}$ and descend to its actual image.
Thus $A^0_{N,L}$ is a unital differential BV algebra.

Any unital differential BV subalgebra containing the current
components contains their $Q$-images and derived brackets.
It therefore contains $C_{N,L}$ and all finite products defining
$A^0_{N,L}$. This proves the asserted minimality within the fixed
ambient distribution algebra. No maximality or topological closure
is asserted.
\end{proof}

\section{Why restriction preserves these products}
\label{sec:hj-bv-restriction}

The ambient restriction $r_N$ already preserves multiplication:
deletion acts independently on each matter slot and commutes with
external tensor product and graded permutations. Thus
\begin{equation}\label{eq:hj-bv-algebra-restriction}
r_N(FG)=(r_NF)(r_NG),\qquad r_N1=1.
\end{equation}
Its failure in \eqref{eq:hj-57} concerns contractions.
Two specific facts about the current space remove that failure
from its generated algebra.

First, in a cross contraction with surviving output $a$, input $b$,
and internal label $c$, a coefficient has the form
$K^F_{ac}H_LK^G_{cb}$, or its reverse orientation. The two
nondecreasing matrices impose
\[
d(b)\leq d(c)\leq d(a).
\]
Low surviving labels therefore force a low contracted label.
This is the bracket compatibility already proved in
\eqref{eq:hj-61}. Second, every current-space element has zero
self-contraction by \eqref{eq:hj-bv-current-closure}.
Nondecreasing degree alone would not imply this second fact.

\begin{proposition}\label{prop:hj-bv-restriction}
The restrictions induce surjective morphisms of unital differential
BV algebras
\[
r_N:A^0_{N+1,L}\longrightarrow A^0_{N,L}.
\]
\end{proposition}
\begin{proof}
The current restriction is a surjection
$C_{N+1,L}\to C_{N,L}$ preserving $Q$ and the bracket, by
Theorem~\ref{thm:hj-current-limit}.
Equation~\eqref{eq:hj-bv-algebra-restriction} therefore maps the
upper generated algebra into the lower one.
Lift a finite polynomial by choosing lifts of its finitely many
current-space factors. This proves surjectivity.

Apply \eqref{eq:hj-bv-all-products} to a product of upper
current-space elements. Every pair bracket is preserved by
$r_N$, and restriction introduces no permutation of the surviving
factors. Consequently,
\begin{equation}\label{eq:hj-bv-delta-restriction}
\boxed{r_N\Delta_L=\Delta_Lr_N\quad\hbox{on }A^0_{N+1,L}.}
\end{equation}
Commutation with $Q$ follows from its ambient restriction identity.
This argument respects actual relations: an upper expression which
is zero as a distribution has zero ambient restriction.
It is not necessary to choose polynomial representatives consistently.
\end{proof}

Set
\begin{equation}\label{eq:hj-bv-hbar}
A_{N,L}=A^0_{N,L}[[\hbar]],\qquad |\hbar|=0,
\end{equation}
with the formal series taken degreewise in the cohomological grading.
For each fixed cohomological degree, coefficients lie in that
degree of $A^0_{N,L}$. The graded algebra has finite sums of
homogeneous cohomological components; no completion in that grading
is imposed. All operations extend coefficientwise, using finite
Cauchy sums for products. A formal series lifts by lifting each
coefficient, so
\begin{equation}\label{eq:hj-bv-formal-restriction}
\boxed{r_N:A_{N+1,L}\twoheadrightarrow A_{N,L}}
\end{equation}
is a surjective morphism of unital differential BV algebras over
$\mathbb C[[\hbar]]$.

\section{A nonzero contraction of actual current products}
\label{sec:hj-bv-nonzero-product}

The restricted algebra has a nontrivial quantum operator, despite
the vanishing of every complete scalar wheel. Put
$D_L=Q+\hbar\Delta_L$. Its square is zero. For homogeneous
$F,G\in A_{N,L}$, its product defect is
\begin{equation}\label{eq:hj-bv-quantum-product}
\begin{aligned}
D_L(FG)-(D_LF)G-(-1)^{|F|}FD_LG
   =\hbar(-1)^{|F|}\{F,G\}_L.
\end{aligned}
\end{equation}
This follows from the derivation rule for $Q$ and
\eqref{eq:hj-bv-bracket}. Thus $D_L$ is not in general a
derivation of multiplication.

To exhibit a nonzero defect in the actual current algebra, take
$N=1$. Choose disjoint open balls $U,V\subset X$ and nonzero,
nonnegative compact smooth functions $\chi_U,\chi_V$ supported
in them. Use the form-degree-zero sources
\[
\alpha=\chi_Ue,\qquad \eta=\chi_Vf,
\qquad F=J(\alpha),\qquad G=J(\eta).
\]
Both currents belong to $C_{1,L}$, have cohomological degree
minus one, and are $\Delta_L$-closed.
The fixed-scale contraction formula is
\begin{equation}\label{eq:hj-bv-current-bracket}
\{F,G\}_L
 =-K_{H_L}(\alpha,\eta)+K_{H_L}(\eta,\alpha),
\end{equation}
where the kernel in \eqref{eq:hj-45} at $n=2$ uses $H_L$ in
place of $P_{0,L}$. This is the positive-scale current bracket
derived in Chapter~\ref{chap:quantum-current-anomaly}.

Select only the external components
\[
\beta_x(z)=b\,\theta_U(z),\qquad
\gamma_x(w)=g\,\theta_V(w)\,
                     d\bar w_1\wedge d\bar w_2,
\]
where $|b|=1$, $|g|=-2$. Choose nonnegative compact
$\theta_U,\theta_V$ with
$\chi_U\theta_U\not\equiv0$ and $\chi_V\theta_V\not\equiv0$,
supported in $U,V$, respectively. All other matter components
are zero. The reversed kernel vanishes because its $\beta$-slot
lies in $V$. The first kernel survives because $(ef)_{xx}=1$.

The surviving heat component is
$d\bar z_1d\bar z_2d\bar w_1d\bar w_2$.
Each orientation identity contributes a factor $4$, so the
coefficient of $bg$ in that first kernel is
\begin{equation}\label{eq:hj-bv-positive-integral}
\frac1{4\pi^2L^2}\int_{X^2}
e^{-|z-w|^2/(4L)}
\chi_U(z)\theta_U(z)\chi_V(w)\theta_V(w)
\,dV_4(z)dV_4(w)>0.
\end{equation}
The integrand is nonnegative, and the two overlap conditions make
it positive on a set of positive measure. The constant is
$4\cdot4/(64\pi^2L^2)$, with $\Omega_z,\Omega_w$ on the right.
Since $|F|=-1$, \eqref{eq:hj-bv-pair} and
\eqref{eq:hj-bv-current-bracket} give
\begin{equation}\label{eq:hj-bv-nonzero-product}
\boxed{\Delta_L(FG)
 =K_{H_L}(\alpha,\eta)-K_{H_L}(\eta,\alpha)\ne0.}
\end{equation}
Thus the algebra generated by the actual currents has a nonzero
BV operator. Spatial separation does not make multiplication a
quantum chain map at a positive heat scale.

\section{The algebraic BV limit and its distribution topology}
\label{sec:hj-bv-limit}

The restriction identities allow all algebra operations to pass
to compatible finite-level coefficients. Define, degreewise,
\begin{equation}\label{eq:hj-bv-limit}
\boxed{A_{\infty,L}=\varprojlim_N A_{N,L}.}
\end{equation}
In a fixed cohomological degree, an element is a compatible array
\[
a=(a_{N,k})_{N\geq1,\ k\geq0},\qquad
a_{N,k}\in A^0_{N,L},\qquad r_Na_{N+1,k}=a_{N,k}.
\]
Each coefficient is a finite polynomial in current-space elements.
There is no uniform polynomial-degree or compact-support bound
across $N$ or $k$. In particular, no completion in observable
arity is taken at a fixed coefficient.

The operations are explicitly
\begin{equation}\label{eq:hj-bv-limit-operations}
\begin{split}
(ab)_{N,k}&=\sum_{i+j=k}a_{N,i}b_{N,j},\\
(Qa)_{N,k}&=Qa_{N,k},\qquad
(\Delta_La)_{N,k}=\Delta_La_{N,k},\\
(D_La)_{N,k}&=Qa_{N,k}+\Delta_La_{N,k-1},\qquad a_{N,-1}=0.
\end{split}
\end{equation}
The unit is the array of unit constant terms and zero positive
$\hbar$ coefficients. Every expression at a specified coordinate
is finite. All BV identities hold after each coordinate
projection, and those projections jointly distinguish elements.
This constructs a unital differential BV algebra, without
identifying it with observables on a single infinite-dimensional
matter space.

A precise distribution topology can be retained with this
algebraic construction. Give each observable-arity summand of
$\mathcal O_{\mathrm{dist},N}$ its strong distribution topology,
and give the algebraic sum over arities the locally convex
direct-sum topology. Give $A^0_{N,L}$ the induced subspace topology.
Then give $A_{N,L}$ the product topology in its $\hbar$ coefficients,
degreewise, and $A_{\infty,L}$ the subspace topology in
$\prod_NA_{N,L}$.

Explicit defining seminorms on the inverse limit have the form
\begin{equation}\label{eq:hj-bv-seminorms}
a\longmapsto
\sum_{m\geq0}\sup_{\phi\in B_m}
\left|\left\langle(a_{N,k})_m,\phi\right\rangle\right|.
\end{equation}
Here $N,k$ are fixed and each $B_m$ is a bounded family of smooth
tests in the corresponding graded symmetric test space.
The sum is finite on each $a_{N,k}$ because that coefficient has
finite observable arity. Finite sums of these seminorms give
the projective topology.

The maps $r_N,Q,\Delta_L$ are continuous linear maps for these
topologies. Differentiation, finite projections of internal labels,
and insertion of the fixed smooth heat kernel send bounded test
families to bounded test families. Transposition therefore gives
continuity on each strong distribution space. Aritywise continuity
extends to the locally convex direct sum, then to coefficient
products and the inverse limit.

For multiplication, if distributions $T,S$ have fixed compact
supports $K,K'$ and orders $p,q$, their external product obeys
\[
|(T\boxtimes S)(\phi)|\leq C_TC_S
\max_{\substack{|\mu|\leq p\\|\nu|\leq q}}
\sup_{K\times K'}|\partial_z^\mu\partial_w^\nu\phi|.
\]
Inserting a fixed smooth heat kernel gives the corresponding
finite-order bounds for brackets. These are fixed-support,
fixed-order estimates. No unrestricted joint continuity or
completeness of $A^0_{N,L}$ is asserted.

\section{Current maps into the BV algebra}
\label{sec:hj-bv-current-inclusion}

The algebra construction changes the target available for products,
but leaves every current component unchanged. The inclusion
\[
S_{N,L}\hookrightarrow A_{N,L}[-1],\qquad uF\longmapsto uF,
\]
is a differential graded Lie map. Its target differential is
$-uD_L$ and its bracket is $u\{F,G\}_L$.
On current-space elements $\Delta_LF=0$, so the differential
agrees with the existing $-uQ$.
The compatibility identity \eqref{eq:hj-64} therefore gives
\begin{equation}\label{eq:hj-bv-current-map}
\boxed{\mathfrak a_+\longrightarrow
S_{\infty,L}\hookrightarrow A_{\infty,L}[-1].}
\end{equation}
It uses exactly the components \eqref{eq:hj-67}.
Every higher identity holds coordinatewise as before.

The ambient counterexample remains: for a removed normalized
pair $b,g$, one has $r_Nb=r_Ng=0$ but
$r_N\{b,g\}_L=1$. In particular, the diagonal quadratic $bg$
is nondecreasing in internal degree, yet its Laplacian is not
preserved. The two conditions used here are bracket compatibility
on the current space and vanishing self-contractions there.
Neither condition should be dropped.

The resulting BV operations are the prescribed signed sums of
contractions. Their compatibility does not assert compatibility
of each individual graph contraction, ordinary matrix traces,
singular propagator operations on arbitrary distributions, or
Gaussian integration over discarded modes. Nor does the algebraic
inverse limit construct factorization maps or an analytic completion.
The trace and geometric comparisons can now be assessed separately
from this established algebra of current products.
